% Abstract: 250 words maximum (BRM). One number per sentence. No receipts tags needed beyond the body's.
Behavioral researchers now prompt large language models with demographic personas and use the answers in place of survey respondents, and the check these simulated samples usually receive is a reading of individual answers.
A simulated sample supports only the uses whose matching validity checks it has passed.
The validity ladder is a protocol of a gate and four rungs.
The gate asks how many repeated draws make a persona's score dependable.
The rungs ask whether a single draw is an answer pattern a person could give, whether gaps between groups match the population's, whether the level matches, and whether the items hang together as they do in people.
Each rung has a statistic, a pass rule, and a use that a pass supports.
We demonstrate the ladder on 28,800 PHQ-8 assessments from four models, read against the National Health and Nutrition Examination Survey.
Persona means become dependable with repeated draws, and every model fails all four rungs above the gate.
Single draws are less varied than real respondents at the same total score.
The gap between low and high income is steepened in every model.
Every model's population mean sits above the survey's after post-stratification.
At level 4, two models have no general factor and the other two fail invariance across income.
Pseudo-models built from survey respondents pass levels 1 and 3 and the level-4 structure rules, and level 2 detects planted distortion but cannot certify faithfulness.
The code, the corpus and a receipt for every reported number are public.
